\section{The original pulled powers have their true full-degree numerical classes}
\label{sec:native-pulled-projective-numerical-classes}

\begin{theorem}[Actual original powers and their genuine numerical classes]
\label{thm:native-pulled-projective-numerical-powers}
\lean{EnvelopeProperty.NativePulledProjectiveNumericalClasses.linePowerIso,EnvelopeProperty.NativePulledProjectiveNumericalClasses.lineClass_power,EnvelopeProperty.NativePulledProjectiveNumericalClasses.classRay,EnvelopeProperty.NativePulledProjectiveNumericalClasses.classRay_apply,EnvelopeProperty.NativePulledProjectiveNumericalClasses.curveDegree_classRay}
\leanok
\uses{thm:native-pulled-projective-line-frames,lem:positive-power-projective-twist,lem:bundle-tensor-power-pullback,lem:native-numerical-line-tensor-power}
\source{bj22-trivval:sec:tc}
The actual original $d$th tensor power of the same original pulled first
twist is genuinely isomorphic to the whole original pulled $O(d)$ line
used above. Its class in the true real numerical quotient is exactly $d$
times the actual original first-twist class. The real class ray is the
literal linear map sending a real coefficient to that multiple of the
same actual original line class. Every actual original closed-curve
functional evaluates this ray by the true full Euler degree of that same
original pulled line on that same original curve.
\end{theorem}
\begin{proof}
The genuine original pullback/tensor-power isomorphism moves the actual
whole first-twist power inside the same original projective pullback.
The genuine original projective twist-power isomorphism identifies it
with the actual $d$th twist, which is the same original pole-sheaf power
used in the original norms and section constructions. Genuine isomorphism
invariance and the actual tensor-power formula in the original numerical
quotient give the stated numerical equality. The scalar class map is
linear by the original vector-space operations. Linearity of every true
closed-curve functional and the genuine full-Euler line-class formula
compute its value on every real coefficient.
\end{proof}

\begin{theorem}[Actual closed curves detect the original very-ample line ray]
\label{thm:native-pulled-projective-numerical-ray}
\lean{EnvelopeProperty.NativePulledProjectiveNumericalClasses.firstLine_isVeryAmple,EnvelopeProperty.NativePulledProjectiveNumericalClasses.curveDegree_firstLine_pos,EnvelopeProperty.NativePulledProjectiveNumericalClasses.classRay_injective_of_curve}
\leanok
\uses{thm:native-pulled-projective-numerical-powers,thm:fc-original-geometric-ample-curve-degree-positive}
\source{bj22-trivval:sec:tc}
If the original closed embedding has the actual original structure-morphism
triangle, its first pulled line is genuinely very ample. On every actual
original closed integral curve its true full-Euler degree is strictly
positive. Hence any actual original closed integral curve makes the real
class-ray map injective. The existence of such a curve is not assumed or
concluded merely from a numerical presentation.
\end{theorem}
\begin{proof}
The given original closed embedding and its literal original triangle
are the actual very-ampleness datum, with the same whole first pulled
line. Genuine very ampleness yields ampleness. The actual geometric
ample restriction theorem gives strictly positive true degree on each
original closed integral curve. Apply that curve functional to equality
of two points on the original class ray. The equality becomes equality
of the two real coefficients times the same strictly positive original
degree, so cancellation gives equality of the coefficients. No numerical
positivity, nonzero-class or injectivity certificate is supplied.
\end{proof}
